% Part: sets-functions-relations
% Chapter: infinite
% Section: dedekind
%
\documentclass[../../../include/open-logic-section]{subfiles}

\begin{document}

\olfileid{sfr}{infinite}{dedekind}
\olsection{Dedekindalgebra's}

We verlangen niet alleen dat de natuurlijke getallen oneindig zijn;
ze moeten ook bepaalde (algebraïsche) eigenschappen hebben: optelling,
vermenigvuldiging enzovoort moeten zich goed gedragen.

Dedekinds idee was om de \emph{opvolgerfunctie} als uitgangspunt te
nemen en de getallen vervolgens te karakteriseren aan de hand van de
volgende eigenschappen:
\begin{enumerate}
	\item Er is een getal, $0$, dat van geen enkel getal de opvolger is
	\\d.w.z., $0 \notin \ran{s}$
	\\d.w.z., $\forall x\ s(x) \neq 0$
	\item Verschillende getallen hebben verschillende opvolgers
	\\d.w.z., $s$ is injectief
	\\d.w.z., $\forall x \forall y (s(x) = s(y) \lif x = y)$
	\item\ollabel{repeatedapplication} Elk getal wordt verkregen door de
	opvolgerfunctie herhaaldelijk op $0$ toe te passen.
\end{enumerate}
De eerste twee voorwaarden zijn eenvoudig met eerste-ordelogica te
behandelen (zie hierboven). Voor \olref{repeatedapplication} volstaat
eerste-ordelogica echter niet. Dedekinds doorbraak was om voorwaarde
\olref{repeatedapplication} als volgt in verzamelingenleer te
herformuleren:
\begin{enumerate}
	\item[3$'$.] De natuurlijke getallen vormen de kleinste verzameling
	die \emph{gesloten is onder de opvolgerfunctie}: als we $s$ op een
	willekeurig element van de verzameling toepassen, verkrijgen we
	opnieuw een element van die verzameling.
\end{enumerate}
We moeten dit echter stap voor stap uitwerken.

\begin{defn}\ollabel{Closure}
	Voor elke functie $f$ heet de verzameling $X$ $f$-\emph{gesloten}
	{dan en slechts dan als} $(\forall x \in X)f(x) \in X$. Definieer
	nu voor elke $o$:
	$$\closureofunder{f}{o} = \bigcap\Setabs{X}{o \in X\text{ en }X
\text{ $f$-gesloten is}}$$
\end{defn}

Dus $\closureofunder{f}{o}$ is de doorsnede van alle $f$-gesloten
verzamelingen die $o$ als element bevatten. Intuïtief is
$\closureofunder{f}{o}$ dan de \emph{kleinste} $f$-gesloten verzameling
die $o$ als element bevat. Het volgende resultaat maakt deze
intuïtieve gedachte precies:
\begin{lem}\ollabel{closureproperties}
	Voor elke verzameling $A$, elke functie $f \colon A \to A$ en elke
	$o \in A$:
	\begin{enumerate}
		\item\ollabel{closurehaselem} $o \in \closureofunder{f}{o}$; en
		\item\ollabel{closureclosed} $\closureofunder{f}{o}$ is $f$-gesloten; en
		\item\ollabel{closuresmallest} als $X$ $f$-gesloten is en $o \in
		X$, dan $\closureofunder{f}{o} \subseteq X$
	\end{enumerate}
\end{lem}

\begin{proof}
Er bestaat ten minste één $f$-gesloten verzameling die $o$ als element
bevat, namelijk $\ran{f}\cup \{o\}$. De verzameling
$\closureofunder{f}{o}$, de doorsnede van \emph{alle} zulke
verzamelingen, bestaat dus. Nu moeten we
\olref{closurehaselem}--\olref{closuresmallest} verifiëren.

\olref{closurehaselem} volgt omdat $o \in \closureofunder{f}{o}$: dit
is een doorsnede van verzamelingen die $o$ allemaal als element hebben.

Voor \olref{closureclosed}, stel dat $x \in \closureofunder{f}{o}$.
Als $o \in X$ en $X$ $f$-gesloten is, dan is dus $x \in X$ en ook
$f(x) \in X$, omdat $X$ $f$-gesloten is. Bijgevolg is
$f(x) \in \closureofunder{f}{o}$.

\olref{closuresmallest} volgt uit het algemene feit dat als $X \in C$,
dan $\bigcap C \subseteq X$.
\end{proof}

Hiermee kunnen we de volgende definitie geven:

\begin{defn}
Een \emph{Dedekindalgebra} bestaat uit een verzameling $A$, een functie
$f \colon A \to A$ en een element $o \in A$ waarvoor geldt:
	\begin{enumerate}
		\item \ollabel{ded:proper} $o \notin \ran{f}$
		\item \ollabel{ded:injection} $f$ is injectief
		\item \ollabel{ded:closure} $A = \closureofunder{f}{o}$
	\end{enumerate}
\end{defn}

Omdat $A = \closureofunder{f}{o}$, volgt uit ons eerdere resultaat dat
$A$ de kleinste $f$-gesloten verzameling is die $o$ als element bevat.
Een Dedekindalgebra is duidelijk Dedekind-oneindig; dit volgt direct
uit de onderdelen \olref{ded:proper} en \olref{ded:injection} van de
definitie. Belangrijker is dat elke Dedekind-oneindige verzameling tot
een Dedekindalgebra kan worden gemaakt.

\begin{thm}\ollabel{thm:DedekindInfiniteAlgebra}
Als er een Dedekind-oneindige verzameling bestaat, dan bestaat er een
Dedekindalgebra.
\end{thm}

\begin{proof}
Zij $D$ Dedekind-oneindig. Dan bestaan er een injectie $g \colon D \to
D$ en een element $o \in D \setminus \ran{g}$. Stel nu $A =
\closureofunder{g}{o}$; volgens \olref{closureproperties} bestaat $A$
en geldt $o \in A$. Stel $f = \funrestrictionto{g}{A}$. We tonen aan
dat $A, f, o$ samen een Dedekindalgebra vormen.

Voor \olref{ded:proper} geldt: $o \notin \ran{g}$ en $\ran{f}
\subseteq \ran{g}$, dus $o\notin \ran{f}$.

Voor \olref{ded:injection}: $g$ is een injectie op $D$ en $f
\subseteq g$, dus is de beperking eveneens een injectie.

Voor \olref{ded:closure}: volgens \olref{closureproperties} is $A$
$g$-gesloten; dus is $A$ zeker ook $f$-gesloten. Daarom is
$\closureofunder{f}{o} \subseteq A$ volgens
\olref{closureproperties}. Ook is $\closureofunder{f}{o}$
$f$-gesloten en geldt $f = \funrestrictionto{g}{A}$; daaruit volgt dat
$\closureofunder{f}{o}$ $g$-gesloten is. Dus
$A \subseteq \closureofunder{f}{o}$ volgens
\olref{closureproperties}.
\end{proof}

\end{document}